% !TeX spellcheck = it_IT

\documentclass[aspectratio=169, ]{beamer}

\usepackage[utf8]{inputenc}
\usepackage[italian]{babel}
\usepackage{hyperref}
\usetheme[titleformat=smallcaps, numbering=fraction, background=light, progressbar=frametitle]{metropolis}

\usepackage{mathtools}
\DeclarePairedDelimiter{\ceil}{\lceil}{\rceil}

\title{Esercitazioni di Informatica B}
\subtitle{Riepilogo MATLAB}

\author{Stefano Cereda\\
	stefano.cereda@polimi.it
}
\date{15/12/2020}
\institute[PoliMi]{Politecnico Milano}
\logo{\includegraphics[width=10mm]{../logopolimi}}

\setbeamercovered{invisible}

\usepackage{xcolor}
\usepackage{listings}

% for matlab code
\definecolor{mygreen}{RGB}{28,172,0} % color values Red, Green, Blue
\definecolor{mylilas}{RGB}{170,55,241}
\definecolor{mGray}{rgb}{0.5,0.5,0.5}
\definecolor{mPurple}{rgb}{0.58,0,0.82}
\definecolor{backgroundColour}{rgb}{0.95,0.95,0.92}

\lstdefinestyle{matlab}{language=Matlab,%
	%basicstyle=\color{red},
	breaklines=true,%
	morekeywords={matlab2tikz},
	keywordstyle=\color{blue},%
	morekeywords=[2]{1}, keywordstyle=[2]{\color{black}},
	identifierstyle=\color{black},%
	stringstyle=\color{mylilas},
	commentstyle=\color{mygreen},%
	showstringspaces=false,%without this there will be a symbol in the places where there is a space
	numbers=left,%
	numberstyle={\tiny \color{black}},% size of the numbers
	numbersep=9pt, % this defines how far the numbers are from the text
	emph=[1]{for,end,break},emphstyle=[1]\color{red}, %some words to emphasise
	%emph=[2]{word1,word2}, emphstyle=[2]{style},
}


\begin{document}
\begin{frame}
    \maketitle
\end{frame}

%\begin{frame}{Acquisizione dati da fogli di calcolo}
%    La funzione \emph{xlsread} consente di aprire in MATLAB file .xls creati con Microsoft Excel.
%
%    Proviamo ad usarla.
%\end{frame}

\begin{frame}{Paghe}
    Data la seguente tabella:

    \begin{tabular}{|c|c|c|c|c|c|}
        \hline
        & op1 & op2 & op3 & op4 & op5 \\
        \hline
        Paga oraria & 5 & 5.5 & 6.5 & 5 & 6.25 \\
        \hline
        Ore settimanali & 40 & 43 & 37 & 50 & 45 \\
        \hline
        Pezzi prodotti & 1000 & 1100 & 1000 & 1200 & 1100 \\
        \hline
    \end{tabular}

    Scrivere uno script Matlab che organizzi la tabella in un array di dati strutturati e risponda alle seguenti domande:
    \begin{enumerate}
        \item Quanto guadagna ogni operaio?
        \item Qual è il salario totale di tutti gli impiegati?
        \item Quanti pezzi vengono prodotti?
        \item Qual è il costo medio di un pezzo?
        \item Quante ore occorrono in media per un pezzo?
        \item Qual è l’operaio più efficiente?
    \end{enumerate}
\end{frame}

\begin{frame}[allowframebreaks]{Archivio mammiferi}

    Si vogliono rappresentare informazioni relative ai mammiferi ospitati in un
    parco protetto.

    In particolare, per ogni mammifero si rappresentano i seguenti dati:\\
    NOME TIPO, TIPO ALIMENTAZIONE, PESO MEDIO, insieme di ESEMPLARI

    Dove ogni esemplare è caratterizzato da:\\
    CODICE ESEMPLARE, ANNO DI NASCITA, NUMERO CUCCIOLI

    \begin{enumerate}
        \item Si crei una variabile mammifero che contenga le informazioni del
            seguente mammifero:\\
            NOME TIPO: Elefante\\
            TIPO ALIMENTAZIONE: Erbivoro\\
            PESO MEDIO: 2500\\
            ESEMPLARE 1: (Codice: 3, Anno: 1985, Cuccioli: 3)\\
            ESEMPLARE 2: (Codice: 5, Anno: 1989, Cuccioli: 0)\\
            ESEMPLARE 3: (Codice: 8, Anno: 1982, Cuccioli: 0)\\

        \item Chieda all'utente se vuole inserire ulteriori elefanti
        \item Si scriva il codice che costruisce, a partire dalla variabile creata al punto
            precedente, l'archivio degli esemplari di elefante sterili sapendo che un
            elefante si dice sterile se è nato prima del 1990 e non ha concepito
            nessun cucciolo.
    \end{enumerate}

    Si stampi su schermo l'elenco dei codici degli elefanti sterili.
\end{frame}

\begin{frame}{Matrice dominante - TdE 20 Luglio 2018}
Una matrice quadrata è definita diagonalmente dominante se il valore assoluto di ciascun elemento sulla
diagonale principale è maggiore della somma dei valori assoluti degli altri elementi sulla stessa riga.

Per esempio, la seguente matrice è diagonalmente dominante:
\begin{tabular}{ccccc}
	20 & 2 & 3 & 4 & 5 \\
	3 & -25 & 5 & 6 & 7 \\
	2 & 3 & 23 & 5 & 6 \\
	8 & 7 & 6 & -30 & 4 \\
	2 & 3 & 4 & 5 & 20 \\
\end{tabular}

Si scriva uno script MATLAB che, ricevuta una matrice $m$, stampi -1 se $m$ non è quadrata, 0 se quadrata ma non diagonalmente dominante, 1 se quadrata e diagonalmente dominante.
\end{frame}

\begin{frame}[allowframebreaks]{Trading - TdE 16 Febbraio 2018}
    Un’agenzia di trading online vuole memorizzare l’andamento del valore dei titoli che controlla.
    La memorizzazione viene effettuata in \textbf{500} istanti temporali equidistanti.
    I dati vengono salvati nel file MATLAB \textbf{log.mat} che contiene:
    \begin{itemize}
        \item la matrice \textbf{titoli}, le cui righe rappresentano i diversi titoli controllati e le cui colonne rappresentano i vari istanti in cui sono stati memorizzati i valori di tali titoli (quindi ogni cella della matrice contiene il valore di un titolo in un dato istante)
        \item il vettore colonna \textbf{andamento}, con lo stesso numero di righe della matrice \textbf{titoli}, che contiene un valore numerico per ogni titolo, indicativo del suo andamento complessivo crescente o decrescente
    \end{itemize}

    \framebreak
    \begin{enumerate}
        \item Scrivere in linguaggio MATLAB una funzione \textbf{splittaMatrice} che:
            \begin{itemize}
                \item riceva in input una matrice \textbf{titoliTot} (con la stessa struttura di \textbf{titoli}), un vettore \textbf{andamentoTot} (con stessa struttura di vettore \textbf{andamento}) e uno scalare \textbf{soglia};
                \item fornisca in output due matrici \textbf{titoliOver} e \textbf{titoliUnder} (ognuna con la stessa struttura di \textbf{titoliTot}). \textbf{titoliOver} include solo le righe di \textbf{titoliTot} corrispondenti agli elementi di \textbf{andamentoTot} con valore maggiore o uguale di soglia. \textbf{titoliUnder}, invece, include le righe di \textbf{titoliTot} corrispondenti agli elementi di \textbf{andamentoTot} con valore minori di soglia.
            \end{itemize}

            \framebreak
        \item Scrivere in linguaggio MATLAB uno script che:
            \begin{enumerate}
                \item legga dal file \textbf{log.mat} i due dati memorizzati: \textbf{titoli} e \textbf{andamento}
                \item richiami la funzione \textbf{splittaMatrice} per separare titoli nelle due matrici \textbf{titoliOver} e \textbf{titoliUnder}, per un valore di soglia pari a 0
                \item crei un vettore \textbf{x} che contenga i 500 istanti di memorizzazione
                \item disegni su due grafici separati (che includano il titolo del grafico e il nome dei due assi) l’andamento dei titoli in \textbf{titoliOver} e \textbf{titoliUnder}, in funzione di x.
            \end{enumerate}
    \end{enumerate}
\end{frame}


\begin{frame}[allowframebreaks]{Cornici - TdE 29 Gennaio 2018}
    Si consideri il seguente problema: si vuole creare una matrice quadrata che sia organizzata come quella in figura.

    \begin{tabular}{|c|c|c|c|c|}
        \hline
        12 & 12 & 12 & 12 & 12 \\
        \hline
        12 & 13 & 13 & 13 & 12 \\
        \hline
        12 & 13 & 14 & 13 & 12 \\
        \hline
        12 & 13 & 13 & 13 & 12 \\
        \hline
        12 & 12 & 12 & 12 & 12 \\
        \hline
    \end{tabular}

    \begin{enumerate}
        \item Si scriva in linguaggio MATLAB una funzione \alert{iterativa} \textbf{cornici} che, data la dimensione N della matrice e un numero di partenza P, restituisca al chiamante una matrice quadrata NxN così definita: la matrice contiene nella cornice più esterna il numero P e numeri crescenti nelle cornici più interne.

        \item Si scriva inoltre uno script in linguaggio MATLAB che acquisisca da tastiera la dimensione desiderata N e il numero di partenza P, invochi la funzione \textbf{cornici} con gli opportuni parametri e infine stampi a video la matrice risultante.

        %\item Si implementi la funzione \alert{ricorsiva} \textbf{corniciRic} in modo che presenti lo stesso comportamento della funzione \textbf{cornici}
        \item Nel TdE c'è un terzo punto che riguarda le funzioni ricorsive
    \end{enumerate}
\end{frame}

\iffalse
\begin{frame}[fragile]{Matlab e binario}
    \framebreak
    Scrivere  una  funzione  Matlab  controllaCP2  che  prende  in  ingresso  una  stringa  (ad
    esempio ‘0101001’)  corrispondente  ad  un  numero  codificato  in  CP2  e  restituisce  due
    valori:
    \begin{enumerate}
        \item \emph{par}  che  vale  true  se  il  numero  è  pari  e  false  altrimenti
        \item \emph{pos}  che  vale  true  se  il  numero  è  positivo  e  false  altrimenti
    \end{enumerate}
    (L’esecuzione  della  funzione  nel  caso  dell’esempio  darebbe  par = false  e  pos =
    true)


    \framebreak

    La funzione matlab deve semplicemente controllare il bit a sinistra per sapere se il numero è positivo, ed il bit a destra per sapere se è pari (ragionate sul perché un numero positivo pari mantiene il primo bit=0 quando passo al negativo).
    \begin{lstlisting}[style=matlab]
function [par, pos] = controllaCP2(stringa)
if stringa(1) == '0'
    pos = true;
else
    pos = false;
end
% versione equivalente:
% pos = stringa(1) == '0';
% par = stringa(end) == '0';
% end
    \end{lstlisting}


\end{frame}



\begin{frame}[allowframebreaks]{Matlab - Ricorsione}
    Un metodo per calcolare il valore approssimato della radice quadrata di un
    numero reale non negativo z, detto metodo babilonese, utilizza la seguente
    relazione ricorsiva:
    \[x(n) = \frac{1}{2}(x(n-1)+\frac{z}{x(n-1)})\]
    \[x(0)=1\]

    Dove $n$ è un intero non-negativo arbitrario, mentre $x(n)$ rappresenta un'approssimazione della radice quadrata del numero $z$. Per esempio, se $z$ è pari a 2 e $n$ è pari a 3 il valore di $x(n)$ ottenuto applicando la formula indicata sopra è 1.4142 (che è una buona approssimazione della radice di 2).

    \framebreak

    L'errore di approssimazione associato a $x(n)$ è definito come segue:
    \[errore(n) = | x(n) - x(n-1)|\]
    e decresce al crescere di $n$. Per esempio, se si calcola la radice di 2 con $n$ pari a 1 si ottiene il valore 1.5000 con errore 0.5000. Se invece si calcola la radice dello stesso valore con $n$ pari a 2 si ottiene il valore 1.4167 con errore 0.0833.

    \framebreak

    \begin{enumerate}
        \item Si scriva una funzione ricorsiva \emph{sqrt1(z,n)} per MATLAB/Octave che
            restituisca un'approssimazione della radice quadrata di $z$ e il relativo errore utilizzando il metodo babilonese. Quando $n$ è pari a zero, si assuma un valore dell’errore pari a inf.
        \item Si scriva una funzione \emph{sqrt2(z,err)} per MATLAB/Octave che, utilizzando la funzione \emph{sqrt1}, restituisca un valore approssimato della radice quadrata di $z$ con errore non superiore al valore $err$ fornito come parametro.
    \end{enumerate}

    Nella soluzione di entrambi i quesiti non è permesso utilizzare alcuna funzione di
    libreria di MATLAB/Octave a parte la funzione \emph{abs(x)} per il calcolo del valore
    assoluto di x.
\end{frame}

\begin{frame}{Matlab - Grafici}
    Scrivere una funzione Octave con parametri $a$, $b$ e $p$ che esegue le seguenti operazioni:
    \begin{enumerate}
        \item  calcola i valori $\cos(x^2)$ per valori di x compresi fra a e b, con passo p, e li memorizza in un vettore y
        \item visualizza su un grafico l'andamento di y;
        \item restituisce al chiamante il massimo valore di $\cos(x^2)$ per $a \leq x \leq b$ ed il valore di x corrispondente. Per effettuare il
            calcolo in questione si può utilizzare la funzione max che, quando prende come parametro un vettore, restituisce due
            valori di cui il primo è il valore massimo trovato nel vettore ed il secondo è l’indice della posizione del vettore in cui tale
            valore è contenuto.
    \end{enumerate}
\end{frame}


\begin{frame}{Matlab - Scomposizione in fattori}
    Scrivere una funzione che scomponga in fattori primi un numero.
    Il risultato dovrà composto da due vettori: il primo conterrà le basi dei fattori, il secondo conterrà gli esponenti dei fattori, cioè quante volte essi sono ripetuti.

    Esempio: il numero 100 è uguale a $2^2 * 5^2$, per cui il vettore delle basi sarà [2 5] mentre quello degli esponenti corrispondenti sarà [2 2].

    Un modo semplice per fattorizzare un numero consiste nel dividerlo ripetutamente per tutti i numeri che gli sono inferiori fino ad ottenere 1.
\end{frame}
\fi

\end{document}
